% =====================================================================
% pertemuan-14.tex
% Pertemuan 14 - Symbol Table, Scope Management, dan Type Checking
% Sub-CPMK 1.6 - Praktikum 4
% =====================================================================
\documentclass[aspectratio=169,10pt,t]{beamer}
\usepackage{config/beamerthemeubb}

\UBBmeeting{14}

\begin{document}

\UBBtitleframe{14}{Symbol Table, Scope Management, dan Type Checking}{Praktikum 4}

% ---------------------------------------------------------------------
\begin{frame}{Capaian Pertemuan}
  \textbf{CPMK induk:} CPMK 01 --- Menjelaskan tahapan kompilasi
  (\textit{lexical}, \textit{syntax}, \textit{semantic analysis}).

  \vspace{3pt}
  \begin{block}{Sub-CPMK yang dinilai}
    \small
    \textbf{Sub-CPMK 1.6:} Mengimplementasikan \textit{symbol table}
    bersarang dan pemeriksa tipe dalam bahasa C, dengan berkas kode
    dan keluaran diagnosis sebagai bukti.
  \end{block}

  \vspace{3pt}
  \begin{exampleblock}{Cakupan Praktikum 4}
    \footnotesize
    Praktikum 4 menilai Sub-CPMK 1.5, 1.6, dan 1.7. Artefak 1.5
    (tabel SDD) dari Pertemuan 13 dan 1.7 (diagnostik kesalahan) ikut
    dikumpulkan.
  \end{exampleblock}
\end{frame}

% ---------------------------------------------------------------------
\begin{frame}{Agenda Pertemuan (150 menit)}
  \small
  \begin{tabular}{@{}p{1.4cm}p{12.3cm}@{}}
    \toprule
    \textbf{Menit} & \textbf{Kegiatan} \\
    \midrule
    20 & Briefing Praktikum 4 dan cakupan artefak \\
    25 & Ceramah: struktur tabel simbol dan scope \\
    30 & Praktik: implementasi \texttt{symtab.c} \\
    35 & Praktik: pemeriksa tipe dan traversal AST \\
    30 & Pengujian, diagnostik, dan perbandingan \\
    10 & Laporan dan penutup \\
    \bottomrule
  \end{tabular}

  \vspace{6pt}
  {\scriptsize\color{ubbblue!75}Porsi praktik 95 menit. Setiap
  mahasiswa menyerahkan simbol tabel, type checker, dan keluaran
  diagnosis.}
\end{frame}

% ---------------------------------------------------------------------
\begin{frame}{Posisi Pertemuan pada Rantai Proyek}
  \UBBchain{5}

  \vspace{2pt}
  \begin{itemize}
    \item Pertemuan ini membangun simpul \textbf{symbol table} dan
          pemeriksa tipe.
    \item Tabel simbol menjadi pusat informasi bagi semantik dan
          validator AST.
    \item Praktikum 4 menghasilkan Milestone M4.
  \end{itemize}
\end{frame}

% ---------------------------------------------------------------------
\begin{frame}{Atribut Simbol dan Operasi Tabel}
  \begin{columns}[T,onlytextwidth]
    \begin{column}{0.46\textwidth}
      \textbf{Atribut simbol}
      \begin{itemize}
        \item nama dan panjangnya;
        \item jenis: variabel, fungsi, parameter;
        \item tipe;
        \item scope asal;
        \item lokasi deklarasi;
        \item alamat atau register.
      \end{itemize}
    \end{column}
    \begin{column}{0.50\textwidth}
      \textbf{Operasi dasar}
      \begin{itemize}
        \item \texttt{define}: tambah ke scope saat ini;
        \item \texttt{lookup}: cari dari dalam ke luar;
        \item \texttt{lookup\_current}: cari di scope saat ini;
        \item \texttt{enter\_scope}: dorong scope;
        \item \texttt{leave\_scope}: lepas scope teratas.
      \end{itemize}
    \end{column}
  \end{columns}

  \vspace{4pt}
  {\footnotesize\color{ubbblue!80}Representasi proyek Subset-C:
  daftar tertaut per scope; hash table dapat dipakai untuk program
  besar.}
\end{frame}

% ---------------------------------------------------------------------
\begin{frame}{Tabel Informasi}
  \scriptsize
  \setlength{\tabcolsep}{4pt}
  \begin{tabular}{@{}p{3.0cm}p{9.6cm}@{}}
    \toprule
    \textbf{Tabel} & \textbf{Informasi yang disimpan} \\
    \midrule
    Identifier & Semua identifier, nomor urut, dan atributnya \\
    Array & Dimensi, tipe elemen, batas, dan ukuran \\
    Blok & Variabel dalam blok yang sama \\
    Real & Elemen bernilai real dan representasinya \\
    String & Panjang, isi, dan lokasi string \\
    Display & Blok yang sedang aktif pada scope bersarang \\
    \bottomrule
  \end{tabular}

  \vspace{4pt}
  {\footnotesize\color{ubbblue!80}Tabel display membantu lookup
  mengetahui blok aktif. Proyek Subset-C memakai satu tabel simbol per
  scope; tabel di atas menjadi acuan bila bahasa diperluas.}
\end{frame}

\begin{frame}{Stack Scope dan Resolusi Nama}
  \begin{columns}[T,onlytextwidth]
    \begin{column}{0.44\textwidth}
      \centering
      \begin{tikzpicture}[
          sc/.style={draw=ubbblue!45,fill=ubbblue!8,rounded corners=3pt,
            minimum width=3.8cm,minimum height=0.6cm,font=\scriptsize,
            text=ubbblue},
          ar/.style={-{Stealth[length=1.5mm]},draw=ubbblue!55,thick}
        ]
        \node[sc] (s0) {Scope global};
        \node[sc,below=0.25cm of s0] (s1) {Scope fungsi};
        \node[sc,below=0.25cm of s1] (s2) {Scope blok};
        \draw[ar] (s2) -- (s1);
        \draw[ar] (s1) -- (s0);
      \end{tikzpicture}
      \par {\scriptsize Pencarian dari dalam ke luar}
    \end{column}
    \begin{column}{0.52\textwidth}
      \begin{itemize}
        \item \texttt{lookup} mencari di scope saat ini.
        \item Jika tidak ada, lanjut ke scope induk dan global.
        \item Tidak ditemukan: identifier tidak dideklarasikan.
        \item Shadowing: identifier dalam menutupi identifier luar.
        \item Deklarasi ganda dalam scope yang sama adalah kesalahan.
      \end{itemize}
    \end{column}
  \end{columns}
\end{frame}

% ---------------------------------------------------------------------
\begin{frame}[fragile]{Struktur Tabel Simbol (\file{symtab.h})}
  \begin{columns}[T,onlytextwidth]
    \begin{column}{0.56\textwidth}
      \lstinputlisting[language=C,firstline=15,lastline=27,
        firstnumber=15,
        caption={Struktur \texttt{Symbol} dan \texttt{Scope}},
        label={lst:p14-symh}]{../buku_ajar/code/subsetc/include/symtab.h}
    \end{column}
    \begin{column}{0.40\textwidth}
      \begin{itemize}
        \item \texttt{Symbol} menyimpan nama, tipe, jenis, dan lokasi.
        \item \texttt{next} membentuk daftar per scope.
        \item \texttt{Scope} menyimpan daftar simbol dan pointer induk.
        \item Scope global tidak memiliki induk.
      \end{itemize}
    \end{column}
  \end{columns}
\end{frame}

% ---------------------------------------------------------------------
\begin{frame}[fragile]{Mendorong Scope Baru}
  \begin{columns}[T,onlytextwidth]
    \begin{column}{0.56\textwidth}
      \lstinputlisting[language=C,firstline=28,lastline=37,
        firstnumber=28,
        caption={Fungsi \texttt{symtab\_enter\_scope}},
        label={lst:p14-enter}]{../buku_ajar/code/subsetc/src/symtab.c}
    \end{column}
    \begin{column}{0.40\textwidth}
      \begin{itemize}
        \item Alokasi scope baru dengan \texttt{calloc}.
        \item \texttt{parent} menunjuk scope saat ini.
        \item Scope baru menjadi current.
        \item \texttt{leave\_scope} membebaskan simbol lalu kembali
              ke induk.
      \end{itemize}
    \end{column}
  \end{columns}
\end{frame}

% ---------------------------------------------------------------------
\begin{frame}[fragile]{\texttt{lookup} pada Scope Bersarang}
  \begin{columns}[T,onlytextwidth]
    \begin{column}{0.58\textwidth}
      \lstinputlisting[language=C,firstline=79,lastline=93,
        firstnumber=79,
        caption={Pencarian simbol per scope},
        label={lst:p14-lookup}]{../buku_ajar/code/subsetc/src/symtab.c}
    \end{column}
    \begin{column}{0.38\textwidth}
      \begin{itemize}
        \item Bandingkan panjang nama lebih dahulu.
        \item \texttt{strncmp} hanya jika panjang sama.
        \item Nama tidak harus null-terminated.
        \item \texttt{symtab\_lookup} mengulang untuk setiap induk.
      \end{itemize}
    \end{column}
  \end{columns}
\end{frame}

% ---------------------------------------------------------------------
\begin{frame}[fragile]{Contoh Shadowing}
  \begin{columns}[T,onlytextwidth]
    \begin{column}{0.52\textwidth}
      \begin{lstlisting}[language=C,numbers=none,
        caption={Dua simbol \texttt{x} pada scope berbeda}]
int x = 1;
{
    int x = 2;
    print(x);
}
print(x);
      \end{lstlisting}
    \end{column}
    \begin{column}{0.44\textwidth}
      \begin{itemize}
        \item Di dalam blok, \texttt{x} merujuk ke 2.
        \item Setelah blok, \texttt{x} kembali ke 1.
        \item Tabel menyimpan dua simbol pada scope berbeda.
        \item Deklarasi ganda pada scope yang sama tetap kesalahan.
      \end{itemize}
    \end{column}
  \end{columns}
\end{frame}

% ---------------------------------------------------------------------
\begin{frame}{Deklarasi Ganda}
  \begin{block}{Aturan}
    \small
    Sebelum \texttt{define}, panggil \texttt{lookup\_current}. Jika
    simbol sudah ada, laporkan kesalahan dan jangan menambahkan simbol
    baru.
  \end{block}

  \begin{itemize}
    \item Shadowing melibatkan scope berbeda.
    \item Deklarasi ganda melibatkan scope yang sama.
    \item Pesan harus membedakan keduanya.
    \item Contoh pesan:
          \texttt{baris 1, kolom 8: identifier 'x' sudah
          dideklarasikan}.
  \end{itemize}
\end{frame}

% ---------------------------------------------------------------------
\begin{frame}{Sistem Tipe Subset-C}
  \begin{columns}[T,onlytextwidth]
    \begin{column}{0.46\textwidth}
      \centering
      \begin{tikzpicture}[
          t/.style={draw=ubbblue!45,fill=ubbblue!8,rounded corners=3pt,
            minimum width=1.5cm,minimum height=0.5cm,font=\scriptsize},
          ar/.style={-{Stealth[length=1.4mm]},draw=ubbblue!55}
        ]
        \node[t] (num) at (0,0) {numerik};
        \node[t] (i) at (-1.8,-0.9) {\texttt{int}};
        \node[t] (f) at (0,-0.9) {\texttt{float}};
        \node[t] (c) at (1.8,-0.9) {\texttt{char}};
        \draw[ar] (i) -- (num);
        \draw[ar] (f) -- (num);
        \draw[ar] (c) -- (num);
      \end{tikzpicture}
    \end{column}
    \begin{column}{0.50\textwidth}
      \begin{itemize}
        \item Tipe dasar: \texttt{int}, \texttt{float},
              \texttt{char}, \texttt{void}, \texttt{string},
              \texttt{bool}.
        \item \texttt{error} dipakai untuk propagasi.
        \item Promosi: \texttt{int}/\texttt{char} ke \texttt{float},
              \texttt{char} ke \texttt{int}.
        \item \texttt{float} ke \texttt{int} memerlukan cast.
      \end{itemize}
    \end{column}
  \end{columns}
\end{frame}

% ---------------------------------------------------------------------
\begin{frame}{Aturan Kompatibilitas}
  \begin{columns}[T,onlytextwidth]
    \begin{column}{0.48\textwidth}
      \begin{block}{Assignment}
        \footnotesize
        \begin{itemize}
          \item tipe sama selalu kompatibel;
          \item \texttt{int} dan \texttt{char} dapat dipromosikan ke
                \texttt{float};
          \item \texttt{char} dapat dipromosikan ke \texttt{int};
          \item \texttt{float} ke \texttt{int} tanpa cast: error.
        \end{itemize}
      \end{block}
    \end{column}
    \begin{column}{0.48\textwidth}
      \begin{block}{Aritmetika}
        \footnotesize
        \begin{itemize}
          \item salah satu \texttt{float}: hasil \texttt{float};
          \item keduanya \texttt{int}/\texttt{char}: hasil
                \texttt{int};
          \item operand non-numerik: \texttt{error};
          \item \texttt{error} dipropagasikan tanpa pesan baru.
        \end{itemize}
      \end{block}
    \end{column}
  \end{columns}
\end{frame}

% ---------------------------------------------------------------------
\begin{frame}[fragile]{Pemeriksaan Ekspresi: \texttt{check\_expr}}
  \begin{columns}[T,onlytextwidth]
    \begin{column}{0.60\textwidth}
      \lstinputlisting[language=C,firstline=55,lastline=69,
        firstnumber=55,
        caption={Cabang identifier pada \texttt{check\_expr}},
        label={lst:p14-expr}]{../buku_ajar/code/subsetc/src/typecheck.c}
    \end{column}
    \begin{column}{0.36\textwidth}
      \begin{itemize}
        \item Identifier dicari di tabel simbol.
        \item Tidak ditemukan: pesan dan \texttt{TYPE\_ERROR}.
        \item Ditemukan: kembalikan tipe simbol.
        \item Cabang literal dan operator mengikuti pola sama.
      \end{itemize}
    \end{column}
  \end{columns}
\end{frame}

% ---------------------------------------------------------------------
\begin{frame}[fragile]{Pemeriksaan Blok dan Deklarasi}
  \begin{columns}[T,onlytextwidth]
    \begin{column}{0.60\textwidth}
      \lstinputlisting[language=C,firstline=142,lastline=155,
        firstnumber=142,
        caption={Awal \texttt{check\_block}},
        label={lst:p14-block}]{../buku_ajar/code/subsetc/src/typecheck.c}
    \end{column}
    \begin{column}{0.36\textwidth}
      \begin{itemize}
        \item Dorong scope sebelum memproses isi blok.
        \item Periksa inisialisasi dengan \texttt{check\_expr}.
        \item \texttt{lookup\_current} mendeteksi deklarasi ganda.
        \item Lepas scope setelah blok selesai.
      \end{itemize}
    \end{column}
  \end{columns}
\end{frame}

% ---------------------------------------------------------------------
\begin{frame}{Aturan Konversi Tipe}
  \scriptsize
  \setlength{\tabcolsep}{4pt}
  \begin{tabular}{@{}p{2.4cm}p{2.4cm}p{2.4cm}p{5.4cm}@{}}
    \toprule
    \textbf{Dari} & \textbf{Ke} & \textbf{Implisit?} &
      \textbf{Catatan} \\
    \midrule
    \texttt{int} & \texttt{float} & ya & promosi \\
    \texttt{char} & \texttt{int} & ya & promosi \\
    \texttt{char} & \texttt{float} & ya & promosi \\
    \texttt{float} & \texttt{int} & tidak & memerlukan cast \\
    \texttt{int} & \texttt{char} & tidak & berpotensi kehilangan data \\
    \texttt{void} & tipe lain & tidak & tidak valid \\
    \bottomrule
  \end{tabular}

  \vspace{4pt}
  {\scriptsize\color{ubbblue!75}Subset-C v1.0 belum mendukung cast;
  \texttt{float} ke \texttt{int} dilaporkan sebagai kesalahan.}
\end{frame}

% ---------------------------------------------------------------------
\begin{frame}{Pesan Diagnosis Kesalahan Tipe}
  \scriptsize
  \setlength{\tabcolsep}{4pt}
  \begin{tabular}{@{}p{4.2cm}p{9.0cm}@{}}
    \toprule
    \textbf{Masukan} & \textbf{Pesan} \\
    \midrule
    \texttt{int x = 3.5;} &
      \texttt{baris 1, kolom 9: tidak dapat menetapkan float ke int} \\
    \texttt{int x; x + "halo";} &
      \texttt{baris 1, kolom 12: operator aritmetika memerlukan operand
      numerik} \\
    \texttt{return q;} &
      \texttt{baris 1, kolom 8: identifier 'q' tidak dideklarasikan} \\
    \texttt{int x; int x;} &
      \texttt{baris 1, kolom 8: identifier 'x' sudah dideklarasikan} \\
    \bottomrule
  \end{tabular}

  \vspace{4pt}
  {\scriptsize\color{ubbblue!75}Pesan menyebut baris, nama, dan jenis
  pelanggaran. \texttt{TYPE\_ERROR} meredam pesan turunan.}
\end{frame}

% ---------------------------------------------------------------------
\begin{frame}{Kegiatan Praktikum: M4}
  \begin{block}{Langkah kerja}
    \footnotesize
    \begin{enumerate}
      \setlength{\itemsep}{1pt}
      \item Implementasikan \texttt{define}, \texttt{lookup}, dan
            scope.
      \item Implementasikan aturan tipe.
      \item Traversal AST untuk deklarasi dan ekspresi.
      \item Tangani \texttt{TYPE\_ERROR}.
      \item Uji program valid dan tidak valid.
    \end{enumerate}
  \end{block}

  \vspace{2pt}
  {\footnotesize\color{ubbblue!80}Kumpulkan artefak 1.5 (tabel SDD)
  dari Pertemuan 13 dan 1.7 (diagnostik kesalahan) yang diselesaikan
  pada Pertemuan 15.}
\end{frame}

% ---------------------------------------------------------------------
\begin{frame}{Rubrik Praktikum 4}
  \small
  \setlength{\tabcolsep}{4pt}
  \begin{tabular}{@{}p{3.4cm}p{8.6cm}p{1.8cm}@{}}
    \toprule
    \textbf{Aspek} & \textbf{Indikator} & \textbf{Bobot} \\
    \midrule
    Kebenaran & Scope, lookup, dan tipe benar & 40\% \\
    Kualitas kode & Struktur, penamaan, komentar, tanpa peringatan &
      25\% \\
    Kelengkapan & Deklarasi, ekspresi, fungsi, diagnostik & 20\% \\
    Pengujian & Kasus normal, shadowing, dan error terdokumentasi &
      15\% \\
    \bottomrule
  \end{tabular}

  \vspace{5pt}
  {\footnotesize\color{ubbblue!80}Rubrik ini sama dengan rubrik artefak
  praktikum pada bagian 10 RPS.}
\end{frame}

% ---------------------------------------------------------------------
\begin{frame}{Tugas dan Artefak}
  \begin{block}{Artefak Sub-CPMK 1.6}
    \textit{Symbol table} bersarang dan pemeriksa tipe dalam bahasa C,
    dengan berkas kode dan keluaran diagnosis sebagai bukti.
  \end{block}

  \small
  \begin{itemize}
    \item Tabel simbol mendukung scope bersarang dan shadowing.
    \item Pemeriksa tipe menangani promosi dan propagasi error.
    \item Minimal 20 kasus uji semantik.
    \item Diagnosis menyebut baris, nama, dan jenis pelanggaran.
    \item Artefak 1.5 (tabel SDD) dan 1.7 (diagnostik) ikut
          dikumpulkan.
    \item Laporan singkat diserahkan.
  \end{itemize}

  \vspace{3pt}
  {\footnotesize\color{ubbblue!75}Praktikum 4 berbobot 12\% dan
  dilaporkan setelah Milestone M4 selesai.}
\end{frame}

% ---------------------------------------------------------------------
\begin{frame}{Rangkuman dan Jembatan}
  \begin{itemize}
    \item Tabel simbol memetakan nama ke tipe, scope, lokasi, dan
          jenis simbol.
    \item Scope bersarang diimplementasikan sebagai stack; lookup
          bergerak dari dalam ke luar.
    \item Shadowing dan deklarasi ganda harus dibedakan.
    \item Pemeriksa tipe memakai aturan promosi dan
          \texttt{TYPE\_ERROR} untuk propagasi.
    \item Diagnosis yang baik menyebut baris, nama, dan jenis
          pelanggaran.
  \end{itemize}

  \vspace{4pt}
  \begin{exampleblock}{Pertemuan berikutnya}
    Pertemuan 15 menutup increment proyek: AST, validator kode, dan
    penanganan kesalahan yang menyebut baris, nama, dan jenis
    pelanggaran.
  \end{exampleblock}
\end{frame}

\end{document}
